% GENERATED FILE. Edit canonical lessons, then run npm run generate:books.
% canonical-source-hash: fnv1a64:6f7f74b24beb6a3c
% canonical-lessons: FR-C442-facade, FR-C442-vitrine, FR-C442-enseigne, FR-C442-gouttiere, FR-C442-halle

\chapter{Façades et vitrines}
\label{ch:fr-facades-et-vitrines}
\hlchaptermodality{442}

\begin{chapteropening}
\textbf{By the end of this chapter:} \emph{I can describe the front of buildings: the façade, the shop window, the shop sign, the gutter and the covered market.}

Complete the last lesson of chapter 442: I can describe the front of buildings: the façade, the shop window, the shop sign, the gutter and the covered market.
\end{chapteropening}

\section[la façade]{la façade — the façade, the front (of a building)}

\label{lesson:FR-C442-facade}



\begin{quote}
Before the new one: say the French for the kiosk, then the French for the cobblestone.
\end{quote}

\subsection*{You'll want to know: la façade}
\textbf{la façade} — \textquotedblleft{}the façade, the front (of a building)\textquotedblright{}.

\begin{quote}
La façade de la mairie est toute blanche.
\end{quote}

\begin{quote}
\emph{The front of the town hall is all white.}
\end{quote}

Figurative: \emph{c'est de la façade} — it is all for show.

\begin{cousinweb}
\emph{façade} was borrowed from Italian \emph{facciata} (EtymDB), built on \emph{faccia} 'face', which goes back to Latin \emph{faciēs}.
\end{cousinweb}

\begin{grammarlens}[title={What you've built}]
One more word for telling, explaining and describing at B1.
\end{grammarlens}

\subsection*{Your turn}
\begin{itemize}
\raggedright
  \item \emph{Say it:} \emph{la façade}
  \item \emph{Say it:} \emph{la façade}, once more
  \item \emph{Say it:} a sentence of your own with \emph{la façade}
  \item \emph{Recall:} say the French for the kiosk, then the French for the cobblestone, then say \emph{la façade} again
\end{itemize}

\begin{tcolorbox}[breakable,colback=teal!4,colframe=teal!35!black,title={Before you move on}]
What does la façade mean? (\textquotedblleft{}The façade, the front (of a building)\textquotedblright{}.) Say it once more.
\end{tcolorbox}
\section[la vitrine]{la vitrine — the shop window; the display case}

\label{lesson:FR-C442-vitrine}



\begin{quote}
Before the new one: say the French for the cobblestone, then the French for the façade.
\end{quote}

\subsection*{You'll want to know: la vitrine}
\textbf{la vitrine} — \textquotedblleft{}the shop window; the display case\textquotedblright{}.

\begin{quote}
Elle regarde les chaussures dans la vitrine.
\end{quote}

\begin{quote}
\emph{She looks at the shoes in the shop window.}
\end{quote}

\emph{Faire du lèche-vitrines} — to go window-shopping.

\begin{cousinweb}
\emph{vitrine} is built on \emph{vitre} 'pane of glass' (EtymDB), which goes back to Latin \emph{vitrum} 'glass'.
\end{cousinweb}

\begin{grammarlens}[title={What you've built}]
One more word for telling, explaining and describing at B1.
\end{grammarlens}

\subsection*{Your turn}
\begin{itemize}
\raggedright
  \item \emph{Say it:} \emph{la vitrine}
  \item \emph{Say it:} \emph{la vitrine}, once more
  \item \emph{Say it:} a sentence of your own with \emph{la vitrine}
  \item \emph{Recall:} say the French for the cobblestone, then the French for the façade, then say \emph{la vitrine} again
\end{itemize}

\begin{tcolorbox}[breakable,colback=teal!4,colframe=teal!35!black,title={Before you move on}]
What does la vitrine mean? (\textquotedblleft{}The shop window; the display case\textquotedblright{}.) Say it once more.
\end{tcolorbox}
\section[l'enseigne]{l'enseigne — the shop sign; the brand}

\label{lesson:FR-C442-enseigne}



\begin{quote}
Before the new one: say the French for the façade, then the French for the shop window.
\end{quote}

\subsection*{You'll want to know: l'enseigne}
\textbf{l'enseigne} — \textquotedblleft{}the shop sign; the brand\textquotedblright{}.

\begin{quote}
L'enseigne du café est en bois.
\end{quote}

\begin{quote}
\emph{The café's sign is made of wood.}
\end{quote}

A feminine noun. \emph{Une grande enseigne} — a big chain of shops.

\begin{cousinweb}
\emph{enseigne} is inherited from Latin \emph{īnsignia} 'marks, badges', from \emph{signum} 'mark' (Lewis \& Short).
\end{cousinweb}

\begin{grammarlens}[title={What you've built}]
One more word for telling, explaining and describing at B1.
\end{grammarlens}

\subsection*{Your turn}
\begin{itemize}
\raggedright
  \item \emph{Say it:} \emph{l'enseigne}
  \item \emph{Say it:} \emph{l'enseigne}, once more
  \item \emph{Say it:} a sentence of your own with \emph{l'enseigne}
  \item \emph{Recall:} say the French for the façade, then the French for the shop window, then say \emph{l'enseigne} again
\end{itemize}

\begin{tcolorbox}[breakable,colback=teal!4,colframe=teal!35!black,title={Before you move on}]
What does l'enseigne mean? (\textquotedblleft{}The shop sign; the brand\textquotedblright{}.) Say it once more.
\end{tcolorbox}
\section[la gouttière]{la gouttière — the roof gutter, the rain gutter}

\label{lesson:FR-C442-gouttiere}



\begin{quote}
Before the new one: say the French for the shop window, then the French for the shop sign.
\end{quote}

\subsection*{You'll want to know: la gouttière}
\textbf{la gouttière} — \textquotedblleft{}the roof gutter, the rain gutter\textquotedblright{}.

\begin{quote}
L'eau de pluie coule dans la gouttière.
\end{quote}

\begin{quote}
\emph{The rainwater runs into the gutter.}
\end{quote}

On the roof; the gutter in the street has another name.

\begin{cousinweb}
\emph{gouttière} is built on \emph{goutte} 'drop', which is inherited from Latin \emph{gutta} (Lewis \& Short).
\end{cousinweb}

\begin{grammarlens}[title={What you've built}]
One more word for telling, explaining and describing at B1.
\end{grammarlens}

\subsection*{Your turn}
\begin{itemize}
\raggedright
  \item \emph{Say it:} \emph{la gouttière}
  \item \emph{Say it:} \emph{la gouttière}, once more
  \item \emph{Say it:} a sentence of your own with \emph{la gouttière}
  \item \emph{Recall:} say the French for the shop window, then the French for the shop sign, then say \emph{la gouttière} again
\end{itemize}

\begin{tcolorbox}[breakable,colback=teal!4,colframe=teal!35!black,title={Before you move on}]
What does la gouttière mean? (\textquotedblleft{}The roof gutter, the rain gutter\textquotedblright{}.) Say it once more.
\end{tcolorbox}
\section[la halle]{la halle — the covered market (hall)}

\label{lesson:FR-C442-halle}



\begin{quote}
Before the new one: say the French for the shop sign, then the French for the roof gutter.
\end{quote}

\subsection*{You'll want to know: la halle}
\textbf{la halle} — \textquotedblleft{}the covered market (hall)\textquotedblright{}.

\begin{quote}
Le samedi, les paysans vendent leurs légumes sous la halle.
\end{quote}

\begin{quote}
\emph{On Saturdays the farmers sell their vegetables in the covered market.}
\end{quote}

\emph{Les halles} — the covered market of a town.

\begin{cousinweb}
\emph{halle} comes from Old French \emph{hale}, from Frankish \emph{halla} 'large room, hall' (EtymDB).
\end{cousinweb}

\begin{grammarlens}[title={What you've built}]
One more word for telling, explaining and describing at B1.
\end{grammarlens}

\subsection*{Your turn}
\begin{itemize}
\raggedright
  \item \emph{Say it:} \emph{la halle}
  \item \emph{Say it:} \emph{la halle}, once more
  \item \emph{Say it:} a sentence of your own with \emph{la halle}
  \item \emph{Recall:} say the French for the shop sign, then the French for the roof gutter, then say \emph{la halle} again
\end{itemize}

\begin{tcolorbox}[breakable,colback=teal!4,colframe=teal!35!black,title={Before you move on}]
What does la halle mean? (\textquotedblleft{}The covered market (hall)\textquotedblright{}.) Say it once more.
\end{tcolorbox}
